\chapter{Обзор проблемы и постановка задач исследования}\label{ch:review}

В данной главе приводится введение в задачу декодирования визуальной информации по сигналам мозга и постановка задач диссертационной работы. Сначала описываются нейробиологические модальности "--- фМРТ и ЭЭГ "--- и их взаимодополняющие пространственно-временные характеристики. Далее дается обзор существующих подходов: одномодальных методов декодирования, контрастивного обучения представлений, генеративных моделей реконструкции и методов анализа пространственно-временной структуры нейросигналов. Затем вводится единый формализм декодирования мультимодальных многомерных временных рядов и формулируются три постановки, отвечающие возрастающей сложности учитываемой структуры данных: оценка гемодинамического лага, пространственно-временная классификация в условиях ограниченной выборки и мультимодальная реконструкция стимулов. Детальное изложение методов и экспериментов приводится в последующих главах, основанных на работах автора~\cite{dorin2024forecasting,dorin2025enhancing,dorin2026decoding,dorin2025mmro,dorin2024conf66mipt,dorin2025conf67mipt,dorin2026conf68mipt,dorin2025aij}.

\section{Нейробиологические модальности и их пространственно-временные характеристики}

%% TODO: фМРТ и BOLD-сигнал, гемодинамический лаг; ЭЭГ и ее временное разрешение;
%% взаимодополняемость модальностей; специфика одновременной регистрации (артефакты МР-среды).
%% Основа: раздел 1 и введение магистерской диссертации.

\section{Обзор существующих методов декодирования}

%% TODO: одномодальные методы (фМРТ, ЭЭГ); контрастивное обучение и CLIP;
%% генеративные модели реконструкции (диффузионные модели, SDXL, IP-Adapter);
%% методы анализа пространственно-временных характеристик (ROI, риманова геометрия);
%% наборы данных с одновременной регистрацией фМРТ--ЭЭГ.
%% Основа: глава 1 магистерской диссертации.

\section{Единая постановка задачи декодирования}

%% TODO: пространства модальностей X_(m), совместное пространство X, пространство стимулов Y,
%% отображение g: X -> Y; три конкретизации (лаг, классификация, реконструкция).
%% Основа: глава 2 магистерской диссертации и тезис ИОИ-2026.

\section{Выводы по главе}

%% TODO: краткие выводы и мотивация следующих глав.
